\documentclass[11pt]{article}
\usepackage[margin=1in]{geometry}
\usepackage[T1]{fontenc}
\usepackage{lmodern,microtype,amsmath,amssymb,amsthm,booktabs,tabularx,enumitem,graphicx,needspace,placeins}
\usepackage[colorlinks=true,linkcolor=blue,urlcolor=blue]{hyperref}
\newtheorem{proposition}{Proposition}
\newtheorem{corollary}[proposition]{Corollary}
\newtheorem{definition}{Definition}
\newcommand{\dd}{\mathrm d}
\newcommand{\Rnet}{\mathcal R}
\title{Housing Allocation, Fertility, and Population\\\large An illustrative model}
\author{Working note for Tommaso De Santo}
\date{September 6, 2026}
\begin{document}
\begin{center}
{\Large Housing Allocation, Fertility, and Population}\par
\smallskip
{\large An illustrative model}\par
\smallskip
{\small Working note for Tommaso De Santo \quad September 6, 2026}
\end{center}
\medskip

Borrowing limits can keep housing with old owners even when young households
would give up more consumption to use it. We show that a small reallocation
can then improve welfare while holding fertility fixed. We also study housing
policy during a demographic transition: an earlier decline in the desire for
children starts the adjustment, and a later policy changes its course.

\section{Environment}

\textbf{Households.}
At date $t$, the economy contains $Y_t$ young households and $O_t$ old
households. A young household has income $y_i>0$, liquid wealth $b_i>0$, and
chooses total nondurable expenditure $c$, total housing $h$, completed
fertility $n>0$, saving $a'\geq0$, and tenure $m\in\{R,O\}$. Fertility is a
continuous choice made once. The pair $(y_i,b_i)$ is drawn from the same
distribution $F$ in each cohort.

\textbf{Preferences.}
Each child uses $\chi>0$ units of goods and $\kappa>0$ units of housing. Define
the parts of the household bundle available to the adults as
\begin{equation}
 x=c-\chi n>0,\qquad s=h-\kappa n>0.
 \label{eq:amend_usable_bundles}
\end{equation}
Young and old flow utilities are
\begin{align}
 u_t^y(c,h,n)&=\log(c-\chi n)+\alpha\log(h-\kappa n)
 +\vartheta_t\log n,
 \label{eq:amend_preferences}\\
 u^o(c^2,h^2,e)&=\log c^2+\gamma\log h^2+\omega_B\log e.
 \label{eq:amend_old_preferences}
\end{align}
Both child requirements reduce the bundle left for adults. All log arguments
and preference weights are positive. The weights $\alpha$ and $\gamma$ value
young and old housing, $\vartheta_t$ values children, and $\omega_B$ values
estates. The factor $\beta>0$ discounts old-age utility. The estate $e$ is valued in goods at
the date after the old period. Parents value what they leave, but their
children do not inherit it as entrant wealth in this model.

Each household draws once an idiosyncratic preference $\xi_i$ for ownership,
independently of $(y_i,b_i)$, and observes it before making its choices.
The draw is logistic with location $\bar\xi$ and scale $\sigma_\xi>0$.

\textbf{Demography.}
A child born at $t$ produces $\nu>0$ young households at $t+1$, after survival
and household formation. Thus
\begin{equation}
 Y_{t+1}=\nu\bar n_tY_t,\qquad O_{t+1}=Y_t,
 \label{eq:amend_demography}
\end{equation}
where $\bar n_t$ is average fertility among the young. Population counts adult households; Appendix~\ref{app:amend_units} gives
the conversion to resident persons.

\textbf{Housing and asset markets.}
The fixed housing stock is $\bar H$. Rental units satisfy
$h\leq h_R^{\max}$, owner units satisfy $h\leq h_O^{\max}$, and
$h_R^{\max}<h_O^{\max}$. The consumption good is the numeraire. Goods and one-period bonds trade
with the rest of the world. The bond price is $q\in(0,1)$ and its gross
return is $R_f=q^{-1}$; housing clears against the domestic stock. The owner
of record pays property tax at rate $\tau^p$. A competitive rental intermediary buys a unit for $P_t>0$ at the
beginning of date $t$, then collects $r_t$, pays $\tau^pP_t$, and holds a unit
worth $P_{t+1}$ at the end of the period. Free entry therefore imposes
\begin{equation}
 r_t+P_{t+1}=R_fP_t+\tau^pP_t
 \quad\Longleftrightarrow\quad
 r_t=(R_f+\tau^p)P_t-P_{t+1}.
 \label{eq:amend_usercost}
\end{equation}
Rent is paid in date-$(t+1)$ goods, so its date-$t$ value is $q r_t>0$;
at a steady state, $r=(R_f-1+\tau^p)P$. There is no capital-gains tax. Property-tax revenue is rebated equally to
all young and old households. The date-$t$ value of this rebate is $T_t$.

\textbf{Home finance.}
A young buyer finances the share $\phi_t\in(0,1)$ and posts
$(1-\phi_t)P_th$ at purchase. Closing occurs before current income and rebates
are available, and unsecured borrowing is unavailable. Thus only liquid
wealth $b_i$ can satisfy the down payment. The owner cannot buy more housing
when old or rent its retained housing to another household. At death the
estate sells it back into the fixed stock, and the proceeds enter warm-glow
estate expenditure.

\section{Household problems}

Households take prices and transfers as given. Renters carry financial wealth
into old age; owners carry the house and its mortgage.

\textbf{Young renters.}
Conditional on renting, the young household solves
\begin{align}
 W_t^R(i)=\max_{c,h,n,a'}\;&u_t^y(c,h,n)+\beta V_{t+1}^R(R_fa') \nonumber\\
 \text{s.t.}\quad&
 c+a'+q\,r_th=y_i+b_i+T_t,\nonumber\\[-2pt]
 &h\leq h_R^{\max},\qquad a'\geq0.
 \label{eq:amend_young_renter}
\end{align}
The renter enters old age with financial wealth $R_fa'$.

\textbf{Young owners.}
Conditional on owning, the young household solves
\begin{align}
 W_t^O(i)=\max_{c,h,n,a'}\;&u_t^y(c,h,n)+\beta V_{t+1}^O(a_{t+1},h) \nonumber\\
 \text{s.t.}\quad&
 c+a'+[(1-\phi_t)+q\tau^p]P_th=y_i+b_i+T_t,\nonumber\\[-2pt]
 &a_{t+1}=R_fa'-R_f\phi_tP_th,\qquad
 (1-\phi_t)P_th\leq b_i,\nonumber\\[-2pt]
 &h\leq h_O^{\max},\qquad a'\geq0.
 \label{eq:amend_young_owner}
\end{align}
The property tax is set aside in the bond. The mortgage is repaid when the
household becomes old, leaving financial wealth $a_{t+1}$ and title $H=h$.

\textbf{Old households.}
An old household has financial wealth $a$. An owner also carries housing $H$
from young age. Renters and owners solve, respectively,
\begin{align}
 V_t^R(a)&=\max_{c^2,h^2,e}u^o(c^2,h^2,e),\nonumber\\
 &c^2+qe+q\,r_th^2=a+T_t,\quad 0<h^2\leq h_R^{\max},
 \label{eq:amend_old_renter}\\
 V_t^O(a,H)&=\max_{c^2,h^2,e}u^o(c^2,h^2,e),\nonumber\\
 &c^2+qe+q\,r_th^2=a+P_tH+T_t,\quad
 0<h^2\leq H,\quad e\geq P_{t+1}h^2.
 \label{eq:amend_old_budget}
\end{align}
The owner may sell part of its house, and retained housing must be included
in its estate. The budget includes sale receipts, property tax, and the value
of housing eventually sold by the estate.

\textbf{Tenure choice.}
A household owns when $W_t^O(i)+\xi_i\geq W_t^R(i)$. The logistic ownership
taste gives the owner share
\begin{equation}
 \pi_t^O(i)=
 \frac{\exp\{[W_t^O(i)+\bar\xi]/\sigma_\xi\}}
 {\exp\{W_t^R(i)/\sigma_\xi\}+\exp\{[W_t^O(i)+\bar\xi]/\sigma_\xi\}}.
 \label{eq:amend_tenure}
\end{equation}
The taste draw smooths aggregate tenure shares without changing conditional
renter or owner policies.

\section{Equilibrium}

Prices and transfers equate housing demand to the fixed stock and balance the
government budget. Let $\bar h_t^Y$ and $\bar h_t^O$ denote average housing per
young and old household, averaging over types and tenure choices. Let $G_t$
be the old cohort's normalized distribution over net financial wealth,
inherited housing titles, and tenure, with separate mass $O_t$.
Housing clearing and the government budget are
\begin{equation}
 Y_t\bar h_t^Y+O_t\bar h_t^O=\bar H,\qquad
 (Y_t+O_t)T_t=q\tau^p P_t\bar H.
 \label{eq:amend_clearing}
\end{equation}
The second condition rebates all property-tax receipts in current value.
Goods and bonds trade externally; their markets do not impose an additional
domestic clearing equation.

\begin{definition}[Equilibrium]
Given policies, fertility preferences, the initial state $(Y_0,O_0,G_0)$,
and the fixed entrant distribution $F$, an equilibrium consists of price and
transfer sequences $(P_t,r_t,T_t)$, household choices, tenure probabilities,
and cohort states $(Y_t,O_t,G_t)$ satisfying:
\begin{enumerate}[label=(\roman*),nosep,leftmargin=*]
 \item Young and old households solve their stated problems, and tenure
 follows \eqref{eq:amend_tenure}.
 \item Rental entry satisfies \eqref{eq:amend_usercost}.
 \item Housing clears and the government budget balances as in
 \eqref{eq:amend_clearing}.
 \item Cohort sizes follow \eqref{eq:amend_demography}; $G_{t+1}$ is the
 distribution of assets, titles, and tenure chosen by the young at $t$.
\end{enumerate}
After an unexpected reform, old households keep their inherited assets and
titles and reoptimize using the new price path. Existing contractual
obligations remain fixed.
\end{definition}

\section{Housing allocation}

We ask whether the planner can improve the use of the existing housing stock
by moving space from old owners to young households. The planner chooses
housing and consumption directly. It respects the physical limits on each
tenure and can relax private financing restrictions. Fertility, tenure,
estates, and all future real allocations remain fixed.

Let $i$ be a young owner and $j$ an old owner. Their marginal values of
housing, measured in units of current consumption, are:
\begin{equation}
 MV_i^Y=\frac{\alpha(c_i-\chi n_i)}{h_i-\kappa n_i}
       =\frac{\alpha x_i}{s_i},\qquad
 MV_j^O=\frac{\gamma c_j^2}{h_j^2}.
 \label{eq:amend_marginal_values}
\end{equation}
These values measure the consumption each household would give up for a
little more housing. Comparing them requires no interpersonal comparison of
utility levels.

\begin{proposition}[Housing misallocation]
\label{prop:amend_welfare}
Suppose matched groups of young and old owners have equal positive mass
and room for a small transfer of housing. Assume uniformly positive young
saving, adult consumption, and room below the young physical housing cap.
If $MV_i^Y-MV_j^O$ is uniformly positive on matched pairs, a sufficiently small reallocation toward the young can
make them better off and leave every other household equally well off, in
the absence of real reallocation costs.
\end{proposition}

\begin{proof}
Give the young owner $\epsilon$ units of the old owner's housing. Keeping
the old owner's utility fixed requires additional consumption:
\begin{equation}
 D_j(\epsilon)=c_j^2\left[
 \left(\frac{h_j^2}{h_j^2-\epsilon}\right)^\gamma-1\right].
 \label{eq:amend_direct_compensation}
\end{equation}
Take this amount from the young owner. Goods and housing balance. Since
$D_j'(0)=MV_j^O$, the young owner's utility gain has derivative
$(MV_i^Y-MV_j^O)/x_i>0$ at zero. A small change preserves feasibility.
The seller replaces the lost estate title with a bond, and the buyer's
later sale of the extra title repays its financing. These claims offset,
so future allocations and existing creditors are unchanged.
Appendix~\ref{app:amend_settlement} gives the payments.
\end{proof}

\textbf{Why the values can differ.}
With positive saving and interior old-age choices, a young owner who would
buy more housing if its down-payment limit were relaxed has $MV_i^Y>q r_t$,
provided the physical owner cap is slack. An old owner with slack retention and estate
bounds instead satisfies:
\begin{equation}
 MV_j^O=\frac{\gamma c_j^2}{h_j^2}=q r_t.
 \label{eq:amend_oldmrs}
\end{equation}
Young buyers need cash at purchase; old owners already hold their homes.
The resulting value gap gives the reallocation in Proposition~\ref{prop:amend_welfare}.
The argument applies at any date when these conditions hold, including
during a demographic transition.



The rental size limit restricts an alternative way to obtain space. A
household can be unable to rent a larger home and unable to fund its
purchase. The owner-to-owner comparison above holds tenure fixed; a renter
already at its physical cap cannot receive more housing in that tenure.
An ownership taste can make constrained ownership preferable even when
larger rentals are available.

If moving $\epsilon$ housing uses $C(\epsilon)$ goods, with $C(0)=0$ and
$C'(0)=K_{ijt}\geq0$, the condition becomes $MV_i^Y>MV_j^O+K_{ijt}$.
The young pays this real cost as well as the old owner's compensation.
A fixed moving cost requires a finite gain large enough to cover it.

The proposition establishes inefficiency relative to a planner who can
relax private finance. It identifies a direction of improvement without
solving the entire first-best allocation. A market-based reform with
households choosing freely requires a separate welfare comparison.

\begin{figure}[!htbp]
 \centering
 \includegraphics[width=.73\textwidth]{output/model/simplified_olg_amendments/theory_slides_misallocation.pdf}
 \caption{Housing moves from the old owner to the young buyer, with the old
 owner compensated in consumption. Each curve uses the household's own
 housing on the horizontal axis. The arrows describe the same transfer of
 space. Fertility and future real allocations are fixed.}
 \label{fig:amend_allocation}
\end{figure}
\FloatBarrier

\section{Housing access and fertility}

More space makes children easier to accommodate. Paying for that space
reduces the goods available to the family. Fertility rises when the first
effect dominates the second.

Consider the young owner's conditional problem at a given date. Assume
positive saving, a binding down payment, a slack physical housing cap, and
interior old-age choices. Suppress household and date indices. Let
$k=\beta(1+\gamma+\omega_B)$ and
$w=y+b+T_t+qT_{t+1}$, lifetime resources before housing. Eliminating saving
from the original problem gives:
\begin{equation}
 (1+k)x+\chi n+u h=w,\qquad u=q r_t.
 \label{eq:amend_owner_resources}
\end{equation}
The net cost of housing is $u$ per unit, after accounting for its eventual
sale. The fertility first-order condition is:
\begin{equation}
 \frac{\vartheta}{n}=\frac{\chi}{x}+\frac{\alpha\kappa}{s},
 \qquad s=h-\kappa n.
 \label{eq:amend_fertility_foc}
\end{equation}
The marginal benefit of a child equals its goods and space requirements,
valued by the household.

Let $p$ be the net lifetime payment for a marginal unit of extra current
housing, including any accompanying transfers. With prices and other
resources fixed, $p=u$; compensation or a subsidy changes that payment.
The budget response is $(1+k)\dd x=-\chi\dd n-p\dd h$.
The implied fertility response is:

\begin{proposition}[A fertility condition]
\label{prop:amend_fertility}
Under these household conditions, more housing increases fertility locally
if and only if:
\begin{equation}
 \frac{\alpha\kappa}{(h-\kappa n)^2}
 >\frac{\chi p}{[1+\beta(1+\gamma+\omega_B)](c-\chi n)^2}.
 \label{eq:amend_fertility_condition}
\end{equation}
The comparison holds at any date, while tenure and the binding constraints
remain unchanged and saving and old-age choices are reoptimized.
\end{proposition}

The condition uses household allocations and the net payment, with no
borrowing multiplier. It is a conditional prediction of the household problem;
a policy's equilibrium effect also includes changes in prices, transfers,
and tenure. Appendix~\ref{app:amend_fertility} gives a sufficient restriction
using only parameters in a stationary specialization, and its extension to
dates with nonincreasing expected house prices. The fixed-fertility welfare
comparison in Proposition~\ref{prop:amend_welfare} does not require this sign.

\section{Fertility decline and policy during the transition}
\label{sec:amend_transition}

The economy need not be in a steady state when policy changes. Consider an
initial steady state with fertility weight $\vartheta_0$ and financed share
$\phi_0$. At date 0 an unexpected permanent fall to
$\vartheta_1=\vartheta_0-\epsilon$, with $\epsilon>0$, starts a baseline
transition. Explaining this initial change in preferences is outside the
exercise.

At a later date $t_p\geq1$, compare continued policy $\phi_0$ with a
permanent increase to $\phi_1=\phi_0+\delta$, where $\delta>0$.
The increase is announced and implemented at $t_p$. The two continuations
therefore start from the same young and old cohorts and the same inherited
assets and housing titles. They face the same $\vartheta_1$ and all other
non-policy primitives. This construction represents a housing policy
introduced during an adjustment already in progress.

The unanticipated policy may cause prices to jump at $t_p$, while inherited
assets and titles remain fixed. With earlier announcement, the two paths
would begin to differ at that date.

\subsection{Steady states and population}
\label{sec:amend_population}

For constant prices and policies, the household problems determine mean
fertility and housing at any candidate $(P,T)$. The old distribution comes
from repeating the young household choices. A positive steady state solves:
\begin{equation}
 \bar n(P,T;\vartheta,\phi)=\frac1\nu,\qquad
 T=\frac{q\tau^pP}{2}\,[\bar h^Y(P,T)+\bar h^O(P,T)].
 \label{eq:amend_stationary_system}
\end{equation}
The first condition is replacement; the second rebates property-tax revenue.
Once a positive root exists, housing clearing determines population:
\begin{equation}
 N_{\rm hh}^*=Y^*+O^*=2Y^*
 =\frac{2\bar H}{\bar h^{Y*}+\bar h^{O*}}.
 \label{eq:amend_population_scale}
\end{equation}
Thus two steady states can have the same replacement fertility and different
population levels. With the fixed stock, a larger population uses less
housing per household on average.

In the mixed-tenure economies used below, a small fall in $\vartheta$
lowers stationary price and population. Housing per young and per old
household rises. A small increase in $\phi$ raises stationary price,
population, and young housing; old housing falls by more than young housing
rises. These are local comparisons around an economy with both renters and
owners.

\subsection{The two continuations}

Let superscripts 0 and 1 denote the baseline and policy paths from $t_p$.
Because their initial young cohorts are identical, the demographic law gives:
\begin{equation}
 \frac{Y_{t_p+k}^1}{Y_{t_p+k}^0}
 =\prod_{j=0}^{k-1}
       \frac{\bar n_{t_p+j}^1}{\bar n_{t_p+j}^0},\qquad
 O_{t+1}^a=Y_t^a\quad(a=0,1).
 \label{eq:amend_transition_population}
\end{equation}
Relative cohort size records the accumulated fertility differences.
This identity is useful along the entire transition, without assuming
that either path is stationary.

\begin{proposition}[A decline followed by a housing policy]
\label{prop:amend_combined}
There is an open set of economies with positive child goods costs,
positive rebated property tax, and positive masses of renters and owners
for which the following holds. A sufficiently small permanent fall in
$\vartheta$ produces a converging baseline transition with lower initial
fertility and a smaller terminal population. A sufficiently small permanent
increase in $\phi$, introduced at any later date on that baseline, has a
converging continuation from the same inherited state and satisfies:
\begin{equation}
 \bar n_{t_p}^1>\bar n_{t_p}^0,\qquad
 N_{{\rm hh},1}^*>N_{{\rm hh},0}^*.
 \label{eq:amend_policy_comparison}
\end{equation}
The same constraints continue to bind along these local transitions.
\end{proposition}

Appendix~\ref{app:amend_transition} gives the assumptions and proof. The
policy need not restore the original population or raise fertility above
replacement.

Both terminal fertility levels equal $1/\nu$. With exponential convergence,
\eqref{eq:amend_transition_population} also implies:
\begin{equation}
 \sum_{t=t_p}^{\infty}\log\frac{\bar n_t^1}{\bar n_t^0}
 =\log\frac{N_{{\rm hh},1}^*}{N_{{\rm hh},0}^*}>0.
 \label{eq:amend_cumulative_fertility}
\end{equation}
The cumulative relative fertility effect is positive. Individual dates can
have effects of either sign: the example admits damped oscillations as it
approaches the new steady state.

\begin{figure}[htbp]
 \centering
 \includegraphics[width=\textwidth]{output/model/simplified_olg_amendments/combined_transition_figure.pdf}
 \caption{A fertility decline followed by a housing policy. From the old
 steady state $S_-$, the preference decline at $t=0$ starts the move toward
 $S_0$. The policy changes
 the intervention-date equilibrium from $I^0$ to $I^1$, with the same inherited
 cohorts and claims, and leads to $S_1$. Dashed lines are constant-price
 schedules. Arrows connect selected first-order equilibrium responses and
 do not imply monotone adjustment. Appendix~\ref{app:amend_figure} defines
 the comparison.}
 \label{fig:amend_credit}
\end{figure}

The permanent policy changes the terminal household problem. A temporary
reform that restores all final primitives cannot change a unique stationary
endpoint if both paths converge to it. The population comparison also does
not establish a welfare gain from the credit reform. The compensated
allocation and the competitive policy answer different questions.

\clearpage
\appendix
\section{Household arguments}
\label{app:amend_households}

\subsection{Financial settlement of the reallocation}
\label{app:amend_settlement}

The direct-allocation proof uses consumption and housing as resources.
Prices enter only to check that inherited claims and future allocations can
be respected. Set $u_t=q r_t$ and write the payment to the seller in excess
of the title price as:
\begin{equation}
 L(\epsilon)=D_j(\epsilon)-u_t\epsilon.
 \label{eq:amend_compensation}
\end{equation}
The seller receives $P_t\epsilon+L(\epsilon)$ and releases the tax reserve
$q\tau^pP_t\epsilon$. After consuming $D_j(\epsilon)$, it buys a bond
costing $qP_{t+1}\epsilon$, replacing the title lost from its estate.
The young buyer spends $D_j(\epsilon)$ less on current consumption and
finances the rest of the purchase and compensation. Its net borrowing is:
\begin{equation}
 (P_t+q\tau^pP_t)\epsilon+L(\epsilon)-D_j(\epsilon)
 =qP_{t+1}\epsilon.
 \label{eq:amend_settlement}
\end{equation}
The two new bond positions offset. At $t+1$ the buyer sells the extra title
for $P_{t+1}\epsilon$ and repays the bond. The old seller's estate receives
the same amount through its bond. The timing and quantity of future housing
sales, estate expenditures, total tax receipts, and common rebates are unchanged.
Existing loans retain their promised repayments. Positive saving provides
room to adjust the original financial portfolio; the planner relaxes the
private borrowing limit. Uniform strict margins on matched groups permit
a common sufficiently small transfer.

With a real cost $C(\epsilon)$, the young owner's exact gain is:
\begin{equation}
 \Delta U_i=
 \log\!\left(1-\frac{D_j(\epsilon)+C(\epsilon)}{x_i}\right)
 +\alpha\log\!\left(1+\frac{\epsilon}{s_i}\right).
 \label{eq:amend_exact_gain}
\end{equation}
Its derivative at zero is $(MV_i^Y-MV_j^O-K_{ijt})/x_i$.
If a fixed cost is present, the finite transfer must keep consumption
positive and make this exact gain positive.

\subsection{Fertility and a sufficient parameter restriction}
\label{app:amend_fertility}

For a fixed tenure, positive saving, and unchanged old-age constraints,
let $\widetilde w$ denote lifetime resources net of any fixed old-age rent.
Let $\widehat h$ be the current housing limit: $h_R^{\max}$ for a renter
and $\min\{h_O^{\max},b/[(1-\phi)P_t]\}$ for an owner. The conditional
problem, after eliminating saving and setting $s=h-\kappa n$, is:
\begin{equation}
 \begin{aligned}
 \max_{x,s,n}\quad &(1+k)\log x+\alpha\log s+\vartheta\log n\\
 \text{s.t.}\quad &(1+k)x+u s+(\chi+\kappa u)n=\widetilde w,
 \qquad s+\kappa n\leq\widehat h.
 \end{aligned}
 \label{eq:amend_reduced_problem}
\end{equation}
The goods and space used by each child cost $\chi+\kappa u$.
For interior old owners, $k=\beta(1+\gamma+\omega_B)$ and
$\widetilde w=y+b+T_t+qT_{t+1}$. For capped old renters,
$k=\beta(1+\omega_B)$ and the old rent must be subtracted:
\begin{equation}
 \widetilde w=y+b+T_t+qT_{t+1}-q^2r_{t+1}h_R^{\max}.
 \label{eq:amend_capped_old}
\end{equation}
This reduction requires different formulas if an old owner's retention or
estate bound binds.

Differentiating the fertility condition and using
$(1+k)\dd x=-\chi\dd n-p\dd h$ gives:
\begin{equation}
 \frac{\dd n}{\dd h}
 =\frac{\alpha\kappa/s^2-\chi p/((1+k)x^2)}{\Delta},\qquad
 \Delta=\frac{\vartheta}{n^2}+\frac{\chi^2}{(1+k)x^2}
              +\frac{\alpha\kappa^2}{s^2}>0.
 \label{eq:amend_fertility_sign}
\end{equation}
The numerator separates the benefit of additional space from the cost of
paying for it.

Differentiating the resource and fertility equations, allowing prices,
resources, and the preference weight to change, gives:
\begin{equation}
 \Delta\,\dd n=\frac{\dd\vartheta}{n}
 +\left(\frac{\alpha\kappa}{s^2}-\frac{\chi u}{(1+k)x^2}\right)\dd h
 +\frac{\chi}{(1+k)x^2}(\dd\widetilde w-h\,\dd u).
 \label{eq:amend_dated_fertility}
\end{equation}
This separates preference, space, and resource effects at any date.
Equation~\eqref{eq:amend_fertility_sign} sets the preference weight fixed and
collects the net resource change into the payment $p$.

For a sufficient condition without the borrowing multiplier, write
$MV^Y=\alpha x/s$. If $\kappa u/\chi>\alpha/(1+k)$ and the housing cap
binds, then $(1+k)\kappa MV^Y>\alpha\chi$. It follows that fertility rises
for any $p<MV^Y$. A condition entirely in parameters can imply this inequality.

\begin{proposition}[A sufficient parameter restriction]
\label{prop:amend_primitive}
Suppose property tax and rebates are zero, old-owner choices are interior,
saving is positive, the physical owner cap is slack, and the down payment
is strictly restrictive.
At stationary prices, the following restriction is sufficient for fertility
to rise after a local housing increase with net payment $p<MV^Y$:
\begin{equation}
 \frac{(1-q)b}{(1-\phi)(y+b)}
 \geq\frac{\alpha}{1+\beta(1+\gamma+\omega_B)+\alpha}.
 \label{eq:amend_fully_primitive}
\end{equation}
It is also sufficient at a transition date if the household expects
$P_{t+1}\leq P_t$.
\end{proposition}

\begin{proof}
Let $v=\kappa u/\chi$ and let $a_h=u\widehat h/\widetilde w$ be the
housing expenditure share at the cap. Within the reduced problem, removing
the current cap gives the optimal housing expenditure share:
\begin{equation}
 g(v)=\frac{\alpha+\vartheta v/(1+v)}{1+k+\alpha+\vartheta}.
 \label{eq:amend_unconstrained_share}
\end{equation}
Strict binding gives $a_h<g(v)$. The function is increasing, and
$g(\alpha/(1+k))=\alpha/(1+k+\alpha)$. At stationary prices,
$a_h=(1-q)b/[(1-\phi)(y+b)]$. The restriction therefore implies
$v>\alpha/(1+k)$. If $P_{t+1}\leq P_t$, then
$u/P_t=1-qP_{t+1}/P_t\geq1-q$, so the same lower bound on $a_h$
holds at that date. The preceding sufficient condition gives the result.
\end{proof}

This parameter restriction excludes very low liquid wealth relative to
lifetime resources. More generally, welfare can increase while fertility
falls: the exact sign still depends on both space and payment.
If a common rental cap increases at both ages, it also raises future rent;
under binding caps the appropriate net payment is $p=u_t+qu_{t+1}$.

A lower preference weight weakly lowers an individual household's fertility
when its full price and transfer sequence is fixed, even allowing tenure
choice. To see this, write its objective as $A(z)+\vartheta\log n(z)$ for a
feasible lifetime plan $z$. Adding the two optimality inequalities at
weights $\vartheta_0$ and $\vartheta_1$ gives:
\begin{equation}
 (\vartheta_1-\vartheta_0)(\log n_1-\log n_0)\geq0.
 \label{eq:amend_preference_fertility}
\end{equation}
The feasible set is unchanged. The general-equilibrium sign also requires
the price and tenure responses considered below.

\subsection{Population units}
\label{app:amend_units}

Let $n^L=a_n n$ count literal children. Consistent rescaling uses
$\nu^L=\nu/a_n$, $\chi^L=\chi/a_n$, and $\kappa^L=\kappa/a_n$.
Replacement literal fertility is $a_n/\nu$; the change in $\log n$ is a
choice-independent constant. If young and old households contain
$a_Y$ and $a_O$ adults and children reside with parents in their birth
period, resident population is:
\begin{equation}
 N_{{\rm persons},t}=a_YY_t+a_OO_t+a_n\bar n_tY_t.
 \label{eq:amend_persons}
\end{equation}
At a positive steady state,
$N_{\rm persons}^*/N_{\rm hh}^*=(a_Y+a_O+a_n/\nu)/2$.
Common fixed counting coefficients therefore preserve the terminal
population ranking. Their empirical values are not needed for this
illustrative comparison.

\section{The combined transition}
\label{app:amend_transition}

The proof uses the original equilibrium equations on a region where the
same household constraints remain strictly binding or slack. Young owners
are at their down-payment limits and below their physical caps; young and
old renters are at their rental caps. Young households save positively.
Old owners have slack retention and estate bounds. All consumption,
usable space, fertility, and estate levels are positive. These conditions
hold uniformly over the fixed compact support of entrant types.

\subsection{Equilibrium equations and inherited claims}

Write $u_t=q r_t$, $a=h_R^{\max}$,
$\rho_O=1+\beta(1+\gamma+\omega_B)$, and
$\rho_R=1+\beta(1+\omega_B)$. Housing choices are
$h_i^O=b_i/[(1-\phi)P_t]$ and $h_i^R=a$.
The conditional resource equations and fertility conditions are:
\begin{align}
 \rho_Ox_i^O+\chi n_i^O&=w_i-u_th_i^O,\nonumber\\
 \rho_Rx_i^R+\chi n_i^R&=w_i-u_ta-qu_{t+1}a,\nonumber\\
 \frac{\vartheta}{n_i^m}&=\frac{\chi}{x_i^m}
  +\frac{\alpha\kappa}{h_i^m-\kappa n_i^m},\qquad
 w_i=y_i+b_i+T_t+qT_{t+1}.
 \label{eq:amend_transition_choices}
\end{align}
The future rent affects renter choices and therefore tenure probabilities.
For the cohorts choosing under these prices, old-owner housing is
$\beta\gamma x_i^O/(qu_{t+1})$ and old-renter housing is $a$.

The aggregate state is $Z_t=(P_t,u_t,Y_t,O_t,M_t,R_t)$, where
$M_t/u_t$ and $R_t$ are total old-owner and old-renter housing.
Given this state, rental entry determines $P_{t+1}$ and fiscal balance
determines $T_t$. The two remaining unknowns
$v=(u_{t+1},Y_{t+1})$ solve:
\begin{align}
 F_1(Z_t,v)&=Y_t\int[\pi_i^Oh_i^O+(1-\pi_i^O)a]\,\dd F
              +M_t/u_t+R_t-\bar H=0,\nonumber\\
 F_2(Z_t,v)&=Y_{t+1}-\nu Y_t
              \int[\pi_i^On_i^O+(1-\pi_i^O)n_i^R]\,\dd F=0.
 \label{eq:amend_transition_residuals}
\end{align}
These are housing clearing and the cohort law. The next state is:
\begin{equation}
 G(Z_t,v)=\left(
 P_{t+1},u_{t+1},Y_{t+1},Y_t,
 \frac{\beta\gamma}{q}Y_t\int\pi_i^Ox_i^O\,\dd F,
 aY_t\int(1-\pi_i^O)\,\dd F\right).
 \label{eq:amend_transition_update}
\end{equation}
If $F_v$ is nonsingular, these equations give a smooth local map
$Z_{t+1}=g(Z_t;\vartheta,\phi)$. It is derived from the original equilibrium,
including old-state evolution.

At an unexpected policy date, let $B,A,H$ denote the inherited old-owner
mass, net financial claims, and purchased titles, respectively. Let
$\mathcal I=(Y_{t_p},O_{t_p},B,A,H)$ and
$L=\gamma/(1+\gamma+\omega_B)$. Initial housing demand satisfies:
\begin{equation}
 Y=Y_{t_p},\quad O=O_{t_p},\quad R=a(O_{t_p}-B),\quad
 M=L[A+PH+T(P,Y,O)B].
 \label{eq:amend_actual_boundary}
\end{equation}
The title price and rebate can change, so $M$ is not predetermined wealth.
Financial claims include the original mortgage repayment; the new $\phi$
does not rewrite that debt. The full conditional inherited states are also
retained to check each old household's feasibility.

Previously chosen saving and rent use the forecasts held before the shock.
In particular, rental entry at $t_p-1$ uses the baseline expected price at
$t_p$, and the old cohort's earlier saving condition need not hold at the
surprisingly realized price. Equations~\eqref{eq:amend_transition_choices}--
\eqref{eq:amend_transition_update} apply to the new continuation from $t_p$
onward, with \eqref{eq:amend_actual_boundary} determining the initial old.

\subsection{Local continuation and the two shocks}

Let $Z^*$ be a positive stationary equilibrium satisfying the stated strict
household conditions. Assume $g$ and the four initial restrictions $b$
are smooth, $F_v$ is nonsingular, and $J=g_Z(Z^*)$ has four eigenvalues
strictly inside and two strictly outside the unit circle. Let $B_0=b_Z$
at the initial stationary data. Its restriction to the stable generalized
eigenspace must be invertible. Finally, at this equilibrium the impact
fertility and stationary population derivatives with respect to both
$\vartheta$ and $\phi$ are strictly positive, holding inherited claims fixed
in the impact comparisons.

Here is the local argument. Write $\lambda=(\vartheta,\phi)$ for the two
parameters. Since one is not an eigenvalue of $J$, the stationary state
$Z^*(\lambda)$ is smooth. Center a continuation on it by
$v_k=Z_{t_p+k}-Z^*(\lambda)$. The exact sequence equations are:
\begin{equation}
 b(Z^*(\lambda)+v_0;\mathcal I)=0,\qquad
 v_{k+1}-g(Z^*(\lambda)+v_k;\lambda)+Z^*(\lambda)=0.
 \label{eq:amend_sequence_equations}
\end{equation}
Centering removes the permanent parameter shift from the tail.
Choose $\eta<1$ above the stable spectral radius and use the sequence norm
$\sup_k\eta^{-k}\|v_k\|$. The derivative in $v$ is
$(B_0v_0,\{v_{k+1}-Jv_k\})$. For a forcing $f_k$, its stable and
unstable components are solved by:
\begin{equation}
 v_k^s=J_s^kc+\sum_{j=0}^{k-1}J_s^{k-1-j}f_j^s,
 \qquad
 v_k^u=-\sum_{j=k}^{\infty}J_u^{k-1-j}f_j^u.
 \label{eq:amend_sequence_inverse}
\end{equation}
The boundary determines $c$ uniquely; both sums are bounded in this norm.
Only the unstable block is inverted, so a zero stable root is allowed.
Uniform household margins make the residual continuously differentiable.
The sequence implicit-function theorem therefore gives a smooth,
exponentially converging continuation for nearby parameters and inherited
data. The same inverse on bounded sequences supplies uniqueness among
continuations staying in a sufficiently small common neighborhood.

For the initial preference decline, inherited data are the original
stationary data. Smoothness gives a constant $C$ such that:
\begin{equation}
 \|Z_t^0(\epsilon)-Z^*(\vartheta_0-\epsilon,\phi_0)\|
 \leq C\epsilon\eta^t,\qquad
 \|Z^*(\vartheta_0-\epsilon,\phi_0)-Z^*(\vartheta_0,\phi_0)\|
 \leq C\epsilon.
 \label{eq:amend_uniform_baseline}
\end{equation}
The entire baseline stays uniformly close to the original equilibrium.
The strict preference derivatives give lower initial fertility and a smaller
endpoint for sufficiently small positive $\epsilon$.

Original household budgets reconstruct assets, titles, and tenure
probabilities smoothly on the fixed compact type support. Their deviations
are therefore uniformly of order $\epsilon$, across every date and type.
This controls individual feasibility as well as the aggregate inherited
condition; closeness of aggregate moments alone would not suffice.
At $\delta=0$, local uniqueness identifies the continuation at $t_p$ with
the baseline tail. Continuity preserves the strict policy derivatives for
all these inherited states in one neighborhood. Integrating them from
$\phi_0$ to $\phi_0+\delta$ gives
\eqref{eq:amend_policy_comparison}. The uniform bound permits common positive
upper bounds on $\epsilon$ and $\delta$, independent of the chosen date $t_p$.
This proves Proposition~\ref{prop:amend_combined} once an economy meeting the
strict conditions is supplied.

\subsection{An economy satisfying the conditions}

Use the following primitive values, with homogeneous income and liquid
wealth and a finite ownership-taste scale:
\begin{align}
 q&=\tfrac12,\quad\phi=\tfrac45,\quad b=\tfrac15,\quad
 \alpha=\beta=\omega_B=\tfrac25,\quad\gamma=\tfrac3{10},\nonumber\\
 \chi&=\tfrac3{20},\quad\kappa=\tfrac12,\quad
 \vartheta=\tfrac{141}{400},\quad\nu=2,\quad
 h_R^{\max}=\tfrac14,\quad h_O^{\max}=2,\quad\sigma_\xi=4,\nonumber\\
 \tau^p&=\tfrac{467}{9250},\quad
 y=\tfrac{3521349281}{1677802000},\quad
 \bar H=\tfrac{68104}{68019}.
 \label{eq:amend_example_primitives}
\end{align}
Set $\bar\xi=\sigma_\xi\log(\pi/(1-\pi))-W^O+W^R$ at
$\pi=11/21$ and the allocation below, then hold both taste parameters fixed
during the shocks. This specifies one preference distribution.

There is an exact stationary equilibrium with $P=Y=O=1$,
$u=9717/18500$, $T=3975571/314587875$, and:
\begin{equation}
 (x_O,h_O,n_O)=(1,1,\tfrac34),\qquad
 (x_R,h_R,n_R)=(\tfrac{99}{74},\tfrac14,\tfrac9{40}).
 \label{eq:amend_example_allocation}
\end{equation}
Mean fertility is $1/2$ and $11/21$ of households own. The young owner's
housing value is $16/25$, exceeding the old owner's value $9717/18500$.
All stated household inequalities have strict margins. Thus the allocation
and transition results apply in the same economy.

Differentiating \eqref{eq:amend_transition_choices} gives the two additional
terms needed for the preference shock: $\dd\vartheta/n_i^m$ in the fertility
condition and $\log n_i^m\,\dd\vartheta$ in the conditional value.
Together with \eqref{eq:amend_tenure}, the latter accounts for endogenous
tenure changes. Let $Q_a=g_a$ for $a\in\{\vartheta,\phi\}$; then:
\begin{equation}
 J=G_Z-G_vF_v^{-1}F_Z,\qquad
 Q_a=G_a-G_vF_v^{-1}F_a,\qquad
 Z_a^*=(I-J)^{-1}Q_a.
 \label{eq:amend_example_derivatives}
\end{equation}
These expressions follow directly from the two implicit equations.

Exact root bounds give one zero root, one real stable root in
$(.00319,.00320)$, a nonreal pair with modulus in $(.41,.43)$, and
unstable roots in $(-80.93,-80.92)$ and $(4.148,4.149)$.
The four initial restrictions are independent on the stable space.
Solving the linear recurrence and those restrictions gives:
\begin{center}
\begin{tabular}{@{}lrr@{}}
\toprule
Derivative at the reference equilibrium & $\vartheta$ increases & $\phi$ increases\\
\midrule
Initial mean fertility & $0.84713$ to $0.84714$ & $0.07966$ to $0.07968$\\
Stationary adult households & $4.79728$ to $4.79729$ & $2.33786$ to $2.33788$\\
\bottomrule
\end{tabular}
\end{center}
The bounds are outward enclosures; all four derivatives are positive.
The preference derivative uses $1<\log(10/3)<2$, with a tighter rational
series bound for the displayed intervals. Strict inequalities persist in
an open neighborhood. Uniformly small income and wealth variation on a
fixed compact probability space preserves the argument as well. This does
not impose zero renter mass, zero child goods costs, or zero property tax.

The local argument above does not give a numerical radius. At the reference
economy, a separate proof with verified interval bounds establishes the
transition and sign comparisons for
$0<\epsilon\leq1/20000$ and $0<\delta\leq1/20000$, uniformly over
$t_p\geq1$. The bounds apply to both infinite paths and preserve the stated
household constraints. This is a small range: the financed share rises from
80\% to at most 80.005\%. The proof and reproducible checks are recorded in
\href{https://github.com/tommasodesanto/Fertility_Spring26/blob/main/output/model/simplified_olg_amendments/oracle_transition_math_assessment.md}{the supporting verification note}.
This bound concerns the reference economy; it is not a common numerical
radius for every economy in the open set. Larger shocks
can change which constraints bind or eliminate the positive endpoint.
For example, the fertility condition implies:
\begin{equation}
 n<\frac{\vartheta h}{\kappa(\alpha+\vartheta)}
 \leq\frac{\vartheta h_O^{\max}}{\kappa(\alpha+\vartheta)}.
 \label{eq:amend_endpoint_obstruction}
\end{equation}
If the last bound is at most $1/\nu$, a positive closed steady state is
impossible. The local construction does not cover such a decline.

\subsection{The two-panel comparison}
\label{app:amend_figure}

Figure~\ref{fig:amend_credit} uses the analytical first-order response at the
economy above. Take a preference decline of size $\epsilon$, introduce the
credit change at $t_p=1$, and set $\delta=\epsilon/2$. These choices specify
a direction as $\epsilon$ tends to zero; they do not select a finite
reform size or a historical date. If $z^a_k$ is the first-order response to
a permanent unit increase in parameter $a$ from the reference equilibrium,
the two displayed directions are:
\begin{equation}
 z_t^0=-z_t^{\vartheta},\qquad
 z_t^1=\begin{cases}
 z_t^0,&t<1,\\
 z_t^0+\tfrac12 z_{t-1}^{\phi},&t\geq1.
 \end{cases}
 \label{eq:amend_figure_paths}
\end{equation}
Their interaction is of second order. The policy response at date 1
satisfies the homogeneous version of \eqref{eq:amend_actual_boundary},
so it preserves that date's inherited cohorts and claims.

The dashed lines are tangents to housing and fertility schedules evaluated
at constant prices and equal age-cohort masses. Write $N=N_{\rm hh}$. At a
candidate $(P,N)$, set $T=q\tau^pP\bar H/N$, solve the household problems with repeated prices,
and define:
\begin{equation}
 \mathcal H(P,N;\vartheta,\phi)
 =\frac N2(\bar h^Y+\bar h^O)-\bar H.
 \label{eq:amend_figure_schedules}
\end{equation}
The left panel solves $\mathcal H=0$ locally for price as a function of
population; the right panel evaluates mean fertility on that same schedule.
On the underlying exact schedules, an intersection with replacement
fertility is a stationary equilibrium. The drawn tangent intersections give
its first-order location. Other points on the schedules need not satisfy
the demographic law.

The solid arrows connect selected transition equilibria. During adjustment,
cohort proportions, inherited assets, and expected future prices differ
from the stationary schedules, so a transition point need not lie on one
of their dashed lines. The arrows compare initial, intervention-date, and
limiting states; they do not impose monotone motion between them. Displayed coordinates are first-order changes per unit $\epsilon$ from the
old equilibrium; symbolic labels refer to economic levels. Both axes use
affine scales.

\end{document}
